% Contrato de encargo de tratamiento de datos personales, para firmar con
% cada cooperativa. Es la versión para imprimir de legal/contrato-encargo.md:
% si una cambia, cambia la otra. Lo ámbar entre corchetes lo completan o
% deciden las partes con su abogado.
\documentclass[10.5pt,a4paper]{article}
% Preámbulo común de los documentos legales (legal/*.tex).
% Se compila cada uno con:   cd legal && latexmk -pdf <archivo>.tex
% La paleta es la del producto y la del manual: si cambia una, cambian las dos.
\usepackage[T1]{fontenc}
\usepackage[utf8]{inputenc}
\usepackage[spanish,es-noshorthands,es-tabla]{babel}
\usepackage[scaled=0.92]{helvet}
\renewcommand{\familydefault}{\sfdefault}
\usepackage{amssymb}
% Los que usan los textos y pdflatex no trae
\DeclareUnicodeCharacter{2713}{\checkmark}
\DeclareUnicodeCharacter{22EF}{\ensuremath{\cdots}}
\DeclareUnicodeCharacter{2192}{\ensuremath{\rightarrow}}
\usepackage[a4paper,top=2cm,bottom=2.2cm,left=2.2cm,right=2.2cm,headheight=14pt]{geometry}
\usepackage[table]{xcolor}
\usepackage{microtype}
\usepackage{enumitem}
\usepackage{booktabs}
\usepackage{tabularx}
\usepackage{array}
\usepackage{longtable}
\usepackage{fancyhdr}
\usepackage{titlesec}
\usepackage{lastpage}
\usepackage[most]{tcolorbox}
\usepackage[hidelinks]{hyperref}

\definecolor{tinta}{HTML}{0B2239}
\definecolor{tinta2}{HTML}{2A4864}
\definecolor{tenue}{HTML}{496A86}
\definecolor{marca}{HTML}{1668C4}
\definecolor{ambar}{HTML}{A67300}
\definecolor{rojo}{HTML}{C2001D}
\definecolor{fondo}{HTML}{F4F7FB}
\definecolor{linea}{HTML}{C3D4E4}
\definecolor{ambarfondo}{HTML}{FDF6E4}
\definecolor{ambarborde}{HTML}{E4C88A}
\definecolor{azulfondo}{HTML}{E7F0FB}
\definecolor{azulborde}{HTML}{C9DDF4}

\color{tinta}
\setlength{\parindent}{0pt}
\setlength{\parskip}{5pt plus 2pt}
\linespread{1.08}
\setlist{itemsep=2pt,topsep=3pt}

\titleformat{\section}{\color{tinta}\large\bfseries}{}{0em}{}
\titleformat{\subsection}{\color{tinta2}\normalsize\bfseries}{}{0em}{}
\titlespacing*{\section}{0pt}{12pt}{4pt}
\titlespacing*{\subsection}{0pt}{8pt}{2pt}

\renewcommand{\arraystretch}{1.25}
\arrayrulecolor{linea}

% Lo que hay que completar a mano: una línea para escribir
\newcommand{\llenar}[1][6cm]{\rule[-2pt]{#1}{0.5pt}}
% Lo que tiene que decidir o completar el abogado o las partes
\newcommand{\completar}[1]{{\color{ambar}[#1]}}
% Una casilla para marcar
\newcommand{\casilla}{\fbox{\rule{0pt}{7pt}\rule{7pt}{0pt}}}

\newtcolorbox{nota}[1][]{colback=azulfondo,colframe=azulborde,boxrule=0.6pt,arc=3pt,
  left=8pt,right=8pt,top=5pt,bottom=5pt,fontupper=\small\color{tinta2},#1}
\newtcolorbox{borrador}{colback=ambarfondo,colframe=ambarborde,boxrule=0.6pt,arc=3pt,
  left=8pt,right=8pt,top=5pt,bottom=5pt,fontupper=\small}

% Pie: qué documento es, su versión y la página
\newcommand{\documento}[2]{%
  \pagestyle{fancy}\fancyhf{}%
  \renewcommand{\headrulewidth}{0pt}%
  \fancyfoot[L]{\scriptsize\color{tenue}#1 · #2}%
  \fancyfoot[R]{\scriptsize\color{tenue}Página \thepage\ de \pageref*{LastPage}}%
}

% Firma: una línea y lo que va debajo
\newcommand{\firma}[1]{\begin{minipage}[t]{0.44\textwidth}\vspace{1.6cm}\rule{\linewidth}{0.5pt}\\{\small\color{tinta2}#1}\end{minipage}}

\begin{document}
\documento{Contrato de encargo de tratamiento de datos personales}{borrador del 11/11/2026}

{\LARGE\bfseries Contrato de encargo de tratamiento de datos personales}\par
{\small\color{tenue}Servicio de control de flota «Control de ruta» · Ley N.º 29733 y su Reglamento (D.S. N.º 016-2024-JUS)}

\begin{borrador}
\textbf{Borrador para revisión de un abogado} (ver \texttt{legal/LEEME.md}).
Lo \completar{ámbar entre corchetes} es lo que hay que completar o decidir.
Las cláusulas técnicas describen lo que el sistema hace de verdad al 11/11/2026:
si el sistema cambia, cambian ellas.
\end{borrador}

Conste por el presente documento el contrato de encargo de tratamiento de
datos personales que celebran:

\begin{itemize}[leftmargin=14pt]
\item \textbf{LA COOPERATIVA}: \llenar[6.5cm], RUC \llenar[3cm],\newline
  con domicilio en \llenar[5.4cm], representada por \llenar[4.6cm],\newline
  identificado con DNI \llenar[2.6cm], según poder inscrito en \llenar[4.4cm];\newline
  en calidad de \textbf{titular del banco de datos y responsable del
  tratamiento}.
\item \textbf{EL ENCARGADO}: \llenar[6.8cm], RUC \llenar[3cm],\newline
  con domicilio en \llenar[5.4cm], representado por \llenar[4.6cm];\newline
  que presta el servicio de control de flota «Control de ruta».
\end{itemize}
En los términos siguientes:

\section{Primera — Objeto}
LA COOPERATIVA encarga a EL ENCARGADO el tratamiento de los datos personales
de sus conductores, cobradores y personal de despacho y gerencia, con la única
finalidad de prestarle el servicio de control de flota contratado
\completar{referencia al contrato de servicio}, conforme a la Ley N.º 29733,
su Reglamento (D.S. N.º 016-2024-JUS) y las instrucciones de LA COOPERATIVA.

\section{Segunda — Datos y finalidades}
\begin{longtable}{>{\raggedright}p{0.44\linewidth}>{\raggedright\arraybackslash}p{0.5\linewidth}}
\toprule
\textbf{Datos} & \textbf{Para qué}\\\midrule\endhead
Identidad en el sistema: nombre, alias, usuario, rol, vehículo & Dar acceso y saber quién maneja qué unidad\\
Ubicación durante el turno, incluida la que se emite con la pantalla apagada o la aplicación en segundo plano & Mapa en vivo, separación entre unidades, vueltas, desvíos, horas en ruta, atención de emergencias\\
Ubicación leída sin enviarse en la pantalla de inicio de turno & Detectar que se está manejando sin haber iniciado el turno\\
Mensajes del chat: texto, notas de voz, fotos, emergencias (SOS) & Coordinación de la operación y atención de emergencias\\
Lugares puntuales del turno, con su hora: dónde se declaró en ruta, ausente o fuera (incluido dónde terminó el turno), cortes de señal, paradas, salidas del trazado & Horas en ruta, señal y presencia, «dónde se traba» la ruta\\
Marcas automáticas («GPS sospechoso», «entrada tardía», «salió solo», «fuera de ruta», «reloj adelantado») & Avisos para que una persona de LA COOPERATIVA los revise. \textbf{No producen ninguna decisión por sí solas}: LA COOPERATIVA se compromete a no sancionar a nadie sólo por una marca automática y a explicarla si el titular lo pide\\
Estado del teléfono durante el turno (batería, ahorro de energía, versión de la app), sin almacenarse & Explicar por qué una unidad dejó de verse\\
Registro de actividad: accesos, altas, bajas, cambios de clave, correcciones & Seguridad y trazabilidad de las acciones administrativas\\
\bottomrule
\end{longtable}
EL ENCARGADO \textbf{no} usará los datos para fines propios ni de terceros: ni
publicidad, ni venta, ni estadísticas para otros clientes, ni el
entrenamiento de sistemas. \completar{Si en el futuro se quisieran estadísticas
agregadas y anónimas del servicio, se acordará por escrito aparte.}

\section{Tercera — Instrucciones}
EL ENCARGADO tratará los datos sólo según las instrucciones de LA COOPERATIVA,
que se expresan en este contrato y en la configuración que ella misma hace del
sistema (sus rutas, sus usuarios, sus permisos). Si alguna instrucción le
parece contraria a la ley, lo dirá por escrito antes de cumplirla.

\section{Cuarta — Confidencialidad}
EL ENCARGADO y su personal guardarán secreto sobre los datos, aun después de
terminado el contrato. Sólo accede a ellos el personal que lo necesite para
dar el servicio y soporte, y queda registrado en el sistema.

\section{Quinta — Medidas de seguridad}
EL ENCARGADO mantiene, como mínimo:
\begin{enumerate}[leftmargin=18pt]
\item Cifrado en tránsito de toda comunicación (HTTPS y WebSocket seguro).
\item Contraseñas almacenadas con función de hash irreversible; nadie puede
  verlas, tampoco EL ENCARGADO. Las sesiones se guardan como huella: una copia
  de la base no permite entrar al sistema.
\item Separación lógica de los datos por cooperativa y control de acceso por
  rol: ninguna cooperativa accede a los datos de otra.
\item Registro de actividad de los accesos y las acciones administrativas,
  consultable por LA COOPERATIVA (sección «Actividad» de su panel).
\item Respaldos automáticos verificados y una copia fuera del servidor
  \textbf{cifrada (AES-256) antes de salir}, que el proveedor de almacenamiento
  no puede leer.
\item Plazos de conservación aplicados automáticamente por el sistema
  (cláusula novena).
\item \completar{Segundo factor de autenticación para el panel de administración del operador.}
\end{enumerate}

\section{Sexta — Sub-encargados}
LA COOPERATIVA autoriza a EL ENCARGADO a servirse de los siguientes
proveedores, sólo para prestar el servicio:

\begin{tabularx}{\linewidth}{>{\raggedright}p{4.2cm}X>{\raggedright\arraybackslash}p{3cm}}
\toprule
\textbf{Proveedor} & \textbf{Servicio} & \textbf{Ubicación}\\\midrule
Railway Corporation & Alojamiento del servidor y la base de datos & \completar{región}\\
\completar{Proveedor del almacenamiento} & Copias de respaldo cifradas & \completar{región}\\
Geoapify & Imágenes del mapa (recibe la dirección IP, no la identidad) & \completar{ }\\
\bottomrule
\end{tabularx}

EL ENCARGADO avisará con \completar{30} días de anticipación cualquier cambio
de proveedor. LA COOPERATIVA podrá oponerse por causa justificada. EL
ENCARGADO responde por sus sub-encargados como por sí mismo.

\section{Séptima — Flujo transfronterizo}
Los datos se alojan en servidores ubicados fuera del Perú. EL ENCARGADO
garantiza que los proveedores de la cláusula sexta ofrecen
\completar{garantías adecuadas: describir, p.\,ej. cláusulas contractuales del
proveedor, nivel de protección del país}. LA COOPERATIVA declara haberlo
informado a los titulares (formulario de consentimiento) \completar{y ante la
ANPD, según corresponda}.

\section{Octava — Incidentes de seguridad}
Si EL ENCARGADO toma conocimiento de un incidente que afecte o pueda afectar
los datos de LA COOPERATIVA, se lo comunicará \textbf{dentro de las 24 horas},
por \completar{teléfono y correo a: }, con lo que se sepa en ese momento: qué
pasó, qué datos y cuántas personas, qué se hizo y qué se recomienda. Así LA
COOPERATIVA puede cumplir su plazo de \textbf{48 horas} ante la Autoridad y,
cuando corresponda, ante los titulares. EL ENCARGADO colaborará con la
investigación y con las medidas que se tomen.

\section{Novena — Conservación}
Los datos se conservan por los plazos siguientes, que el sistema aplica solo:

\begin{tabularx}{\linewidth}{>{\raggedright}p{7.6cm}X}
\toprule
\textbf{Dato} & \textbf{Plazo}\\\midrule
Chat (texto, fotos, notas de voz) & 30 días\\
Mensajes marcados como evidencia por LA COOPERATIVA & 365 días\\
Vueltas, tramos y desvíos & 120 días\\
Lugares puntuales del turno y marcas automáticas & 120 días\\
Turnos y emergencias (SOS) & 365 días\\
Registro de actividad & accesos 90 días; operación 365 días; gestión 730 días\\
Posición en vivo y estado del teléfono & sólo el último valor\\
Copias de respaldo & 30 días\\
\bottomrule
\end{tabularx}

\completar{LA COOPERATIVA podrá pedir plazos distintos por escrito; algunos se
configuran por cooperativa y otros para todo el servidor.}

\section{Décima — Derechos de los titulares}
Si un titular le pide a EL ENCARGADO acceso, rectificación, cancelación,
oposición o portabilidad, éste lo derivará a LA COOPERATIVA en \completar{2}
días hábiles y la asistirá para responder en plazo, entregándole lo que el
sistema tenga de esa persona.

\section{Undécima — Fin del contrato}
Al terminar el contrato, EL ENCARGADO:
\begin{enumerate}[leftmargin=18pt]
\item Entregará a LA COOPERATIVA, si lo pide dentro de los \completar{30} días,
  una copia de sus datos en formato abierto (CSV).
\item Borrará los datos de LA COOPERATIVA del sistema dentro de los
  \completar{30} días siguientes. Las copias de respaldo que los contengan se
  borran solas a los 30 días de hechas.
\item Le confirmará el borrado por escrito.
\end{enumerate}
Salvo que una ley le obligue a conservar algo, en cuyo caso lo dirá.

\section{Duodécima — Verificación}
LA COOPERATIVA podrá pedir información sobre el cumplimiento de este contrato,
y EL ENCARGADO la dará por escrito en \completar{15} días hábiles.
\completar{Auditoría presencial: decidir si se acepta y con qué aviso.}

\section{Decimotercera — Responsabilidad}
Cada parte responde por los incumplimientos que le sean imputables. EL
ENCARGADO que use los datos para un fin distinto del encargado será
considerado responsable del tratamiento por ese uso.

\section{Decimocuarta — Vigencia}
Este contrato rige mientras dure la prestación del servicio y sus
obligaciones de confidencialidad, devolución y borrado subsisten después.

\vspace{6pt}
Firmado en \llenar[4cm], el \llenar[4cm], en dos ejemplares.

\begin{samepage}
\noindent\firma{\textbf{LA COOPERATIVA}\\Nombre, cargo y firma}\hfill\firma{\textbf{EL ENCARGADO}\\Nombre, cargo y firma}
\end{samepage}
\end{document}
